% titel: Blatt 0, A7: Bruchgleichungen

% Teil I: Werkzeuge, kein Lösungsweg.

% Grundidee: Bei "Bruch = Bruch" kreuzweise multiplizieren, aber erst die
% Nenner!=0 als Definitionsbereich sichern; am Ende Ergebnisse prüfen.

Bruchgleichungen kombinieren zwei Gefahren: (i) eine \emph{Division} ist
nur definiert, wenn der Nenner $\neq0$ ist, und (ii) man will den Bruch
auflösen, indem man über Kreuz multipliziert --- eine Aktion, die man nur
für zulässige Nenner anwenden darf. Wer zuerst die Definitionslücken
notiert, verfängt sich nicht in der zweiten Falle.

\begin{werkzeugkasten}[Bruchgleichungen: der Ablauf]
\begin{enumerate}
  \item \textbf{Definitionsbereich festlegen}: Nenner müssen $\neq0$ sein,
        sonst ist der Bruch undefiniert. Diese Bedingungen \emph{vorab}
        notieren.
  \item \textbf{Über Kreuz multiplizieren} (bei ``Bruch = Bruch``):
        $\frac{a}{b}=\frac{c}{d}\ \Rightarrow\ ad=bc$ (für $b,d\neq0$).
        Bei mehreren Brüchen alle auf einen gemeinsamen Nenner bringen
        oder geschickt zusammenfassen.
  \item Die entstehende Gleichung lösen.
  \item \textbf{Prüfen}, ob die Lösungen im Definitionsbereich liegen.
\end{enumerate}
\textbf{Warum über Kreuz?} Multipliziert man $\frac ab=\frac cd$ auf
beiden Seiten mit $b\cdot d$, kürzen sich die Nenner und es bleibt
$ad=bc$. Diese eine Multiplikation ``entfesselt`` die Unbekannte aus den
Nennern.
\end{werkzeugkasten}

\begin{werkzeugkasten}[Definitionslücken systematisch]
Aus $x+2\neq0$ folgt $x\neq-2$, aus $2x-3\neq0$ folgt $x\neq\tfrac32$
usw. Lösungen, die einen Nenner verschwinden lassen, sind \emph{keine}
gültigen Lösungen, auch wenn sie die umgeformte Gleichung erfüllen. Das
Notieren der Ausnahmen \emph{vor} dem Umformen ist entscheidend.
\end{werkzeugkasten}

\begin{achtung}[Typische Fehler]
\begin{itemize}
  \item \textbf{Nenner-Bedingungen vergessen:} Bei $\frac1{x+2}=\frac9x$
        sind $x+2\neq0$ und $x\neq0$ vorab auszuschließen.
  \item \textbf{Vorzeichen in $-...$ beim Umformen:}
        $2-x\neq x-2$; ein f\"uhrendes Minus vor einem Bruch
        ($-\frac{2}{2-x}$) multipliziert \emph{jeden} Nennerterme.
  \item Am Ende nicht nur ``gelöst``, sondern auch gesagt, dass die
        Lösung keinen Nenner aufhebt.
\end{itemize}
\end{achtung}
